\SourceLeaf{327}%
take the particular values below under the substitution \(X=T,\ Y=B\)
\[
\begin{aligned}
&f,\ f',\ f'',\ f''',\ f'''',\ f^{\mathrm{v}}\\
&g,\ g',\ g'',\ g''',\ g'''',\ g^{\mathrm{v}}\\
&h,\ h',\ h'',\ h''',\ h'''',\ h^{\mathrm{v}}
\end{aligned}
\]
whereas they take the following values under the substitution \(X=\mathrm{T},\ Y=\mathrm{B}\)
\[
\begin{aligned}
&\mathrm{f},\ \mathrm{f}',\ \mathrm{f}'',\ \mathrm{f}''',\ \mathrm{f}'''',\ \mathrm{f}^{\mathrm{v}}\\
&\mathrm{g},\ \mathrm{g}',\ \mathrm{g}'',\ \mathrm{g}''',\ \mathrm{g}'''',\ \mathrm{g}^{\mathrm{v}}\\
&\mathrm{h},\ \mathrm{h}',\ \mathrm{h}'',\ \mathrm{h}''',\ \mathrm{h}'''',\ \mathrm{h}^{\mathrm{v}}
\end{aligned}
\]
By \textsc{Taylor}’s theorem, the value of \(Z\) for \(u=R-\frac12r\) expands into the series
\[
\mathrm{L}-\tfrac12r\mathord{.}\frac{\mathrm{d}Z}{\mathrm{d}u}+\tfrac18rr\mathord{.}\frac{\mathrm{dd}Z}{\mathrm{d}u^2}-\tfrac1{48}r^3\mathord{.}\frac{\mathrm{d}^3Z}{\mathrm{d}u^3}+\text{etc.}
\]
and its value for \(u=R+\frac12r\) into
\[
\mathrm{L}+\tfrac12r\mathord{.}\frac{\mathrm{d}Z}{\mathrm{d}u}+\tfrac18rr\mathord{.}\frac{\mathrm{dd}Z}{\mathrm{d}u^2}+\tfrac1{48}r^3\mathord{.}\frac{\mathrm{d}^3Z}{\mathrm{d}u^3}+\text{etc.}
\]
where the derivatives must take their particular values corresponding to \(u=R\), namely
\[
\begin{aligned}
\frac{\mathrm{d}Z}{\mathrm{d}u}&=\mathrm{h}\\
\frac{\mathrm{dd}Z}{\mathrm{d}u^2}&=\mathrm{fh}'+\mathrm{gh}''\\
\frac{\mathrm{d}^3Z}{\mathrm{d}u^3}&=\mathrm{ff'h}'+\mathrm{fg'h}''+\mathrm{f''gh}'+\mathrm{gg''h}''+\mathrm{ffh}'''+2\mathrm{fgh}''''+\mathrm{ggh}^{\mathrm{v}}
\end{aligned}
\]
Since these two values of \(Z\) are, respectively, \(=L+\frac12l\) and \(L-\frac12l\), we obtain
\begingroup\fontsize{10}{12}\selectfont
\[
\begin{aligned}
L={}&\mathrm{L}+\tfrac18(\mathrm{fh}'+\mathrm{gh}'')rr\quad .\ .\ .\ .\ .\ .\ .\ .\ .\ .\ .\ .\ .\ .\ .\ .\ .\ .\ .\ .\ .\quad(4)\\
l={}&-\mathrm{h}r-\tfrac1{24}(\mathrm{ff'h}'+\mathrm{fg'h}''+\mathrm{f''gh}'+\mathrm{gg''h}''+\mathrm{ffh}'''+2\mathrm{fgh}''''+\mathrm{ggh}^{\mathrm{v}})r^3\ .\ .\quad(5)
\end{aligned}
\]
\endgroup
where the first equation is accurate up to but excluding quantities of fourth order, and the second up to but excluding quantities of fifth order\footnote{The orders are measured by treating \(\frac ra\) as a quantity of first order. It is easy to see that the coefficients of \(r,\ rr,\ r^3\), and so on, implicitly contain the divisors \(a,\ aa,\ a^3\), and so on.}.
